\section{Proofs of the fixed-node results}
\label{app:permutation-exchangeability}

\subsection{Monte Carlo permutation exchangeability}

The proof of \cref{lem:plusone-superuniform} separates two facts. First,
statistics computed from i.i.d.\ random permutations are exchangeable. Second,
under A0.2 to A0.4, the usual computation with one observed statistic and $B$
random permutations has the same null rank distribution as a construction with
a randomized baseline \citep{HemerikGoeman2018ExactTestingRandomPermutations}.

Fix a node $t$, feature $j$, and a statistic $T(\cdot, \cdot; U)$ taking values
in a measurable space $(\mathcal{S},\mathcal{A})$. The space may be the real
line, as for a Stage~A score or a single threshold score, or a space of
vectors, as for the vector of scores over a candidate set that the max-type rule
standardizes. When $T$ is real-valued, larger values are more extreme. Let
$\Pi_0, \Pi_1, \dots, \Pi_B$ be i.i.d.\ uniform random permutations of
$\{1,\dots,n_t\}$, independent of the data and of $U$, and define
%
\begin{equation}
  T_b := T(X_{t,j}, \Pi_b(Y_t); U) \qquad (b=0,1,\dots,B).
\end{equation}
%
Conditional on $(X_t, Y_t, U)$, each $T_b$ is a measurable function of $\Pi_b$,
and $(\Pi_0,\dots,\Pi_B)$ are i.i.d., so $(T_0, T_1, \dots, T_B)$ is
exchangeable. Averaging over $Y_t$ preserves exchangeability, so the vector is
also exchangeable conditional on $(X_t,U)$.

\begin{lemma}[The observed statistic has the same null rank distribution]
\label{lem:identity-vs-random}
Assume A0.2 to A0.4, and let $T$ take values in the measurable space
$(\mathcal{S},\mathcal{A})$. Let $\Pi_1,\dots,\Pi_B$ be i.i.d.\ uniform random
permutations of $\{1,\dots,n_t\}$, independent of $(X_t,Y_t,U)$, and set
%
\begin{equation}
  T_{\mathrm{obs}} := T(X_{t,j}, Y_t; U),
  \qquad
  T_b := T(X_{t,j}, \Pi_b(Y_t); U) \quad (b=1,\dots,B).
\end{equation}
%
Let $\Pi_0$ be an additional independent uniform random permutation and set
$T_0 := T(X_{t,j}, \Pi_0(Y_t); U)$. Under the null, conditional on $(X_t,U)$,
%
\begin{equation}
  \big(T_{\mathrm{obs}}, T_1,\dots,T_B\big)
  \stackrel{d}{=}
  \big(T_0, T_1,\dots,T_B\big).
\end{equation}
%
In particular, when $T$ is real-valued, the standard right-tail $+1$ $P$-value
$P_{\mathrm{id}} := (1 + \sum_{b=1}^B \ind\{T_b \ge T_{\mathrm{obs}}\})/(B+1)$
and its randomized-baseline counterpart
$P_{\mathrm{rnd}} := (1 + \sum_{b=1}^B \ind\{T_b \ge T_0\})/(B+1)$ have the same
conditional distribution.
\end{lemma}

\begin{proof}
Fix $(X_t,U)$ and work under the null. By A0.2, permuting the node responses
does not change the conditional distribution of $Y_t$. Use the
permutation-action convention under which
$(\Pi_b \circ \Pi_0^{-1})(\Pi_0(Y_t)) = \Pi_b(Y_t)$, and set
$Y_t' := \Pi_0(Y_t)$. Then $Y_t' \stackrel{d}{=} Y_t$ conditional on $(X_t,U)$.

For $b=1,\dots,B$, define $\Pi_b' := \Pi_b \circ \Pi_0^{-1}$, so
$\Pi_b'(Y_t') = \Pi_b(Y_t)$. Conditional on $(Y_t,\Pi_0)$, right composition by
$\Pi_0^{-1}$ maps the independent uniform permutations $\Pi_b$ to independent
uniform permutations $\Pi_b'$, and this conditional law does not depend on the
conditioning values. Because $Y_t'$ is a function of $(Y_t,\Pi_0)$, the
transformed permutations are independent of $(Y_t',\Pi_0)$. Therefore
%
\begin{equation}
  \begin{aligned}
    \big(T_{\mathrm{obs}}, T_1,\dots,T_B\big)
    &= \big(T(X_{t,j},Y_t;U), T(X_{t,j},\Pi_1(Y_t);U),\dots,T(X_{t,j},\Pi_B(Y_t);U)\big) \\
    &\stackrel{d}{=} \big(T(X_{t,j},Y_t';U), T(X_{t,j},\Pi_1'(Y_t');U),\dots,T(X_{t,j},\Pi_B'(Y_t');U)\big) \\
    &= \big(T_0, T_1,\dots,T_B\big),
  \end{aligned}
\end{equation}
%
conditional on $(X_t,U)$. Applying the same comparison map to these
equal-in-distribution vectors gives the claim for $P_{\mathrm{id}}$ and
$P_{\mathrm{rnd}}$. No step uses the order structure of $\mathcal{S}$, so the
identity holds for a statistic in any measurable space.
\end{proof}

\subsection{Proof of \texorpdfstring{\cref{lem:plusone-superuniform}}{Lemma 1}}
\label{app:proof-plusone}

Fix $(X_t,U)$ and one feature $j \in F_t$, and work under the nodewise complete
permutation null. By \cref{lem:identity-vs-random}, it is enough to prove the
rank bound for an exchangeable tuple, which we denote by
$(T_{\mathrm{obs}},T_1,\dots,T_B)$. Define
%
\begin{equation}
  R := 1 + \sum_{b=1}^B \ind\{T_b \ge T_{\mathrm{obs}}\},
  \qquad
  P_{\mathrm{MC}}^{\ge} = \frac{R}{B+1}.
\end{equation}
%
The left-tail version applies the same argument to $-T$.

\paragraph{No ties.}
If the statistics almost surely have no ties, $R$ is the descending rank of
$T_{\mathrm{obs}}$ among $(T_{\mathrm{obs}},T_1,\dots,T_B)$. Exchangeability
gives $\PP(R = r \mid X_t,U) = 1/(B+1)$ for $r=1,\dots,B+1$, so for any
$\alpha \in [0,1]$,
%
\begin{equation}
  \PP(P_{\mathrm{MC}}^{\ge} \le \alpha \mid X_t,U)
  = \PP\!\big(R \le (B+1)\alpha \mid X_t,U\big)
  = \frac{\lfloor (B+1)\alpha \rfloor}{B+1}
  \le \alpha.
\end{equation}

\paragraph{Ties.}
Let $V_{\mathrm{obs}},V_1,\dots,V_B \stackrel{\mathrm{i.i.d.}}{\sim}
\mathrm{Unif}(0,1)$ be independent of all other randomness, and order the pairs
$(T_b,V_b)$ lexicographically. The augmented tuple is exchangeable and almost
surely has no ties, so the no-ties argument gives a superuniform $P$-value
%
\begin{equation}
  \tilde{P}_{\mathrm{MC}}^{\ge}
  :=
  \frac{
    1 + \sum_{b=1}^B
    \ind\{
      T_b > T_{\mathrm{obs}}
      \text{ or } (T_b = T_{\mathrm{obs}} \text{ and } V_b > V_{\mathrm{obs}})
    \}
  }{B+1}.
\end{equation}
%
Each tie-broken indicator is bounded by $\ind\{T_b \ge T_{\mathrm{obs}}\}$, so
$P_{\mathrm{MC}}^{\ge} \ge \tilde{P}_{\mathrm{MC}}^{\ge}$ and
%
\begin{equation}
  \PP(P_{\mathrm{MC}}^{\ge} \le \alpha \mid X_t,U)
  \le \PP(\tilde{P}_{\mathrm{MC}}^{\ge} \le \alpha \mid X_t,U) \le \alpha.
\end{equation}
%
Taking expectations over $(X_t,U)$ gives the unconditional bound. The
corollaries apply this single-test bound with the Bonferroni union bound; no
independence among the tests is required.

\subsection{Proof of \texorpdfstring{\cref{cor:stageB-maxt}}{Corollary 3}}
\label{app:proof-maxt}

The argument reduces to \cref{lem:plusone-superuniform} in four steps.
(1)~Under the complete null, the $B+1$ rows of the impurity matrix
$(S_{b,c})$, computed from the observed response and from $B$ independent
uniform permutations of it, are exchangeable. This is
\cref{lem:identity-vs-random} applied to the vector-valued statistic
$c \mapsto S_{b,c}$ in $\mathbb{R}^{K}$; the lemma never uses that the statistic
is scalar. (2)~Column standardization uses the mean and standard deviation of
all $B+1$ entries of each column, which are symmetric functions of the rows, so
the standardized matrix has exchangeable rows as well. The exclusion of
candidates with an empty child or zero column variance is likewise a function of
the whole column. (3)~The row minimum is one fixed function applied to every
row, so $T_0,\dots,T_B$ are exchangeable. (4)~The inclusive-rank $P$-value of an
exchangeable tuple is superuniform by the proof of
\cref{lem:plusone-superuniform}, applied to $-T$ for the left tail. The
guarantee is weak familywise control over the $K$ candidates, the probability
of declaring any split under the complete nodewise null.

\subsection{Proof of \texorpdfstring{\cref{prop:adaptive-size}}{Proposition 1}}
\label{app:proof-adaptive}

\paragraph{Reduction to a function of the exceedance probability.} Write
$\theta$ for the per-draw exceedance probability given the observed statistic
and the node data. Given $\theta$, the draws are independent, so the exceedance
count after $r$ draws is $\mathrm{Binomial}(r,\theta)$. The stop and report
decisions are deterministic functions of $(r,e)$ at the batch boundaries, so the
rejection probability $g(\theta)$ is a fixed function of $\theta$.

\paragraph{Monotonicity.} The function $g$ is nonincreasing. Flipping one
indicator from zero to one raises $e$ at every later boundary. A futility stop
(posterior mass at least $\gamma$ above $\alpha$) can then only come earlier, and
a significance stop (mass at least $\gamma$ below $\alpha$) only later. At a
significance stop or at the full budget, the reported value $(e+1)/(r+1)$ is at
least as large on the flipped path. The remaining step, that a significance stop
reports a value strictly below $\alpha$ and a futility stop a value at least
$\alpha$, is the hypothesis of the proposition.

\paragraph{From $g$ to statements (i) to (iii).} Under the complete null the
observed statistic is exchangeable with the $N$ orbit values, so without ties
$\theta$ is uniform on $\{1/N,\dots,1\}$, which is (ii). The continuous
idealization (i) replaces that grid by the uniform distribution on $(0,1)$.
Under it the binomial mixture is the P\'olya urn started from one ball of each
colour, whose next indicator is one with probability $(e+1)/(r+2)$, and
$\rho = \int_0^1 g$ by a finite recursion over the stopping states. For
nonincreasing $g$ the grid average $N^{-1}\sum_{i=1}^{N} g(i/N)$ is at most
$\int_0^1 g$ and at least $\int_0^1 g - (g(0)-g(1))/N \ge \rho - 1/N$, which is
the interval in (ii). Counting ties as exceedances raises $\theta$, and given
$\theta$ the same nonincreasing $g$ gives (iii).
